\scanpage{318}

\mspage{57}
% archive sheet number: 317

Pour ne pas faire dépendre la définition du domaine des $\pi_i$ $(i\in[0,r-1])$ des entiers $\ell_i$, on peut aussi \unsure{considérer} les $\pi_i$ comme des applications
\[
\mathbb N^*\overset{\pi}{\unsure{\text{ --- }|\longrightarrow}}\mathbb N
\qquad(\mathbb N^*=\{i\in\mathbb N\mid i\geq1\}
\]
à supports finis (i.e. tels que $\pi(\unsure{n|i})=0$ pour $i$ grand), et exiger de façon précise
\[
\begin{aligned}
&\operatorname{supp}\pi_0\subset\{1\}\\
&\operatorname{supp}\pi_1\subset[1,\operatorname{long}(\pi_0)]\\
&\quad\vdots\\
&\boxed{\operatorname{supp}\pi_i\subset[1,\operatorname{long}(\pi_{i-1})]}
\quad\text{pour}\quad\boxed{i\in[0,r-1]}
\end{aligned}
\tag*{($**$)}
\]

D'autre part, \deleted{$\unsure{f_i}$} posant $\ell_i=\operatorname{long}(\pi_{\unsure{i-1|i}})$ (pour $i\in[1,\unsure{r|r-1}]$)
\[
 f_i:\Delta_\varepsilon^{\ell_{i-1}}\longrightarrow\Delta_\varepsilon^{\ell_i\,(=\operatorname{long}\pi_{i-1})}
\]
est défini comme
\[
 f_i=f_{\varepsilon,\pi_{i-1}(1)}*f_{\varepsilon,\pi_{i-1}(\unsure{2|1})}*
 \cdots*f_{\varepsilon,\pi_{i-1}(\ell_{i-1})},
\]
de sorte que l'élément $(\unsure{\alpha|*|x})$ de $\operatorname{Nerf}_r(\underline{\mathrm{Ty}})$ est déterminé de façon unique par la donnée des \deleted{$\unsure{\pi}$} $\varepsilon\in\{\pm1\}$ et de
\[
\pi_{\bullet}=(\pi_i)_{\deleted{\unsure{0\leq}}i\in[0,r-1]}
 \in\big(\operatorname{Hom}_{\mathrm{Ens}}(\mathbb N^*,\mathbb N)\big)^r
\]
satisfaisant les conditions $(**)$. \deleted{\unsure{D}} Soit $F_{\varepsilon,\pi_\bullet}$ \deleted{\unsure{le|l'élé}} cet élément de $\operatorname{Nerf}_r(\underline{\mathrm{Ty}})$, \deleted{Alors les}\inserted{et} \deleted{$\unsure{(F_{\varepsilon,\pi_\bullet})_{\varepsilon,\pi_\bullet}}$}\inserted{soit} $N_r$ l'ensemble des $\pi_\bullet$ qu'on vient d'expliciter, Alors
\[
(F_{\varepsilon,\pi_\bullet})_{(\varepsilon,\pi_\bullet)\in\{\pm1\}\times N_r}
\]
est un système de générateurs libre du monoïde $\operatorname{Nerf}_r(\underline{\mathrm{Ty}})$,

Si $\ell_i=0$\inserted{(i.e. $\boxed{\pi_{i-1}=0}$)}, on doit avoir \inserted{$\operatorname{supp}\pi_{\unsure{i+1|i-1}}=\varnothing$ i.e.} $\pi_{\unsure{i-1|i+1}}=0$, d'où $\ell_{i+1}=0$, et de $=$ pour les $\ell_j$ et $\pi_j$ suivants. Géométriquement, \deleted{\unsure{cette}}\inserted{l'hyp.} s'exprime par la condition que le but
